\item Evangelista Torricelli fue la primera persona en darse cuenta de que los seres humanos viven en el fondo de un océano de aire. Él conjeturó correctamente que la presión de la atmósfera es atribuible al peso del aire. La densidad del aire a $0^\circ\text{C}$ en la superficie de la Tierra es $1.29 \text{ kg/m}^3$. La densidad disminuye al aumentar la altitud (a medida que la atmósfera se adelgaza). Por otra parte, si se supone que la densidad es constante a $1.29 \text{ kg/m}^3$ hasta cierta altura $h$ y es cero sobre dicha altura, en tal caso $h$ representaría la profundidad del océano de aire. (a) Use este modelo para determinar el valor de $h$ que da una presión de $1.00 \text{ atm}$ en la superficie de la Tierra. (b) ¿La cima del monte Everest se elevaría sobre la superficie de tal atmósfera?
